\section{Estimating \ours}
\label{sec:method}
In this section, we begin by clearly defining the research problem in Section \ref{31}, followed by a detailed framework of estimating \ours\ in Section \ref{32}.

 \begin{figure}[t]
\centering \footnotesize
% \resizebox{\columnwidth}{!}{
%     \includegraphics{figures/ours.png}%
% }\vspace{-.3em}
\includegraphics[width=\linewidth]{sec/figures/method_new.pdf}
\caption{
    An overview of estimating \ourstitle, which mainly consists of three steps: Descriptive Sub-sentence Identification, Atomic Fact Generation, and Fact Verification. These steps are implemented by three modules: Recognizer, Decomposer, and Verifier. The \underline{underlined part} denotes recognized descriptive content.
}\label{fig:ours}
\vspace{-2mm}
\end{figure}

\subsection{Task and Settings}
\label{31}
Suppose we have an image $I$ and a corresponding prompt $Q$. We then feed them into the LVLM denoted as $\mathcal{M}$, to obtain the generated answer $A$. 
% using the following equation: 
% testing dataset, denoted as $\mathcal{D}={(P_1, I_1), \cdots, (P_N, I_N)}$, where each tuple $(P_i, I_i)$ consists of a prompt $P_i$ and an image $I_i$. The total number of samples in the testing set is represented by $N$. 
% We then input the prompt $P_i$ and the corresponding image $I_i$ into a large vision-language model denoted as $\mathcal{M}$, and obtain the generated answer $A_i$ using the following equation:
% \begin{equation}
% A = \mathcal{M} (Q, I).
% \end{equation} 
Our objective is to design a scoring framework to estimate \ours\ $f$ based on the input prompt $Q$, the input image $I$, and the generated answer $A$. 
It is defined as:
% follows:
% \begin{equation}
$s = \mathcal{F} (A, Q, I)$.  
% \end{equation}
% where the output score 
$s$ is a scalar value ranging between 0.0 and 1.0. Notably, the devised evaluation method is reference-free and doesn't require a ground truth answer. 

\subsection{The Evaluation Framework}
\label{32}
In order to estimate \ours of the generated answers, we introduce a novel framework to implement the scoring function $\mathcal{F}$. The framework comprises three key steps: descriptive sub-sentence identification, atomic fact generation, and fact verification, as depicted in Figure~\ref{fig:ours}. 
% We introduce Recognizer, Dcomposer, and Verifier modules, to fulfill these steps, respectively. 
% \xd{this part doesn't sound right after we change \ours\ to ``estimating \ours''}

\paragraph{Descriptive Sub-sentence Identification.} 
% Unlike LLMs, 
Faithfulness in the context of LVLMs refers to the consistency between the input visual content and the generated answer. Notably, we focus on the details in the answer that describe the input image objectively, to obtain a more precise and fine-grained understanding of the hallucination. 
% as shown in Section~\ref{fig:prompt_des} of the Appendix. 
% \xd{give a natural language definition or reference the prompt in the appendix for ``objectively''}. 
% When faithfulness is low, a significant portion of the generated information may be unrelated to the input and can be considered as hallucination. Hence, unlike LLMs, the primary objective of evaluating faithfulness in LVLMs is to ascertain the authenticity of the visual information contained in the generated answer. 
As shown in Figure~\ref{fig:hallucination}, only some sub-sentences (\ie those with the underline) from the answer require verification with the image input.
 Hence, we need to identify the descriptive sub-sentences from the answer using a recognizer. The sub-sentences denote the short sentences that are split by punctuation in the answer. 
% , to obtain a more precise and fine-grained understanding of the hallucination.

Humans are capable of distinguishing descriptive sub-sentences from other contents (referred to as analytical sub-sentences) by analyzing the semantics of the answers generated by LVLMs. However, manually identifying descriptive sub-sentences is a resource-intensive process, requiring plenty of human labor. To address this problem, we turn to ChatGPT to implement the recognizer as a practical solution, since it has demonstrated 
remarkable semantics understanding capabilities across a wide range of natural language processing tasks~\cite{chatgpt}. Section~\ref{sec:42} shows that ChatGPT can achieve promising performance on this task.

To be more specific, our approach first crafts a prompt $P$ that encompasses task instructions and $K_1$ in-context learning examples. We feed this designed prompt along with the to-be-processed answer $A$ into the ChatGPT, generating the recognized results, defined by the equation 
% \begin{equation}
$\hat{A} = ChatGPT(A, P)$, 
% \end{equation}
where $\hat{A}=\{\{a_1,l_1\},\cdots,\{a_k,l_k\}\}$ signifies the generated result, in which the answer is split into sun-sentences $a$, and every sub-sentence is assigned a label $l$ (\ie descriptive or analytical). Then we extract all descriptive sub-sentences denoted as $A'=\{a'_1,\cdots,a'_t\}$. For a more comprehensive understanding of the specific prompt $P$ utilized in this process, please refer to Section~\ref{prompts} of the Appendix.

% \begin{table*}[]
%     \centering
%     \begin{tabular}{l|cccc}
%         \toprule
%          \bf Model& \bf Conversation & \bf Detailed Description &  \bf Complex Question & \bf Overall\\ 
%          \midrule
%          MiniGPT-4 & 0.7246& 0.5840& 0.8489& 0.7192\\
%          LLaVA & 0.7562& 0.6105&0.7323 & 0.6997\\
%          InstructBLIP & 0.7958 & 0.5760 & 0.7779& 0.7165\\
         
%          Multimodal-GPT&0.7638 &0.5631&0.8926&0.7398\\
%          mPLUG-Owl & 0.8044& 0.6432& 0.8601& 0.7692\\
%          \bottomrule
%     \end{tabular}
%     \caption{Evaluation results of different vision-language models on the LLaVA-D dataset.}
%     \label{llavadataset}
% \end{table*}

\paragraph{Atomic Fact Generation.} 
Despite we have identified descriptive sub-sentences from the answer, there are still multiple facts hybrid in each sub-sentence.
Each descriptive sub-sentence consists of multiple pieces of information (\ie atomic facts), each of which may contain hallucination. Therefore, to access a fine-grained evaluation, we design a decomposer to further break the sub-sentences into atomic facts. 
% Specifically, faithful information corresponds to details present in the input visual modality. 
% In order to assess the faithfulness of answers generated by LVLMs at a fine-grained level, prior research~\cite{Li2023EvaluatingOH} has employed the Caption Hallucination Assessment with Image Relevance (CHAIR) metric~\cite{DBLP:conf/emnlp/RohrbachHBDS18} to evaluate hallucination in image captioning tasks.
% However, it is worth noting that the CHAIR metric primarily focuses on coarse-grained object-level hallucination. 
% , neglecting aspects such as object attributes and relationships between objects. Additionally, it relies on the availability of additional labels for objects and encompasses a limited variety of object types. 
% \xd{maybe less details in CHAIR, it's already covered in related work. or use 2-3 sentences ``Instead of only ...''}
% Hence, to address its limitation, we introduce a novel mechanism that can break down the original sentence into atomic facts. 
% In our pursuit of evaluating faithfulness from fine-grained perspectives, inspired by the TIFA metric~\cite{TIFA}, 
In particular, 
% we break sentences down into atomic facts. W
we define atomic facts as an element belonging to an entity, relation, or attribute, inspired by the existing works~\cite{factscore,TIFA}. 
% \xd{to provide 1-2 examples of \textsc{other}}, 
% which can cover all the aspects of the original sentences. 
% The \textsc{other} attribute includes all attributes of the entity which is not contained by the other four categories (\eg ``The building is very tall''). 
% Specifically, we define atomic fact as the information belonging to one category in the above five categories. 
Importantly, the atomic fact is a minimal unit of information. This handling can ensure the verification of \emph{each} element in the answer without being disturbed by other information. 
% For example, for the category \textsc{entity}, each atomic fact does not involve more than two entities. 
Atomic facts include three types: entity existence, attributes, and relations. An entity fact indicates an object’s existence. Attribute facts relate to characteristics like color and shape. Relation facts describe inter-entity relationships, e.g., the spatial relation. 
In Figure~\ref{fig:hallucination}, we show some examples of atomic facts. 

% Refer to the prompt of atomic fact generation in Section~\ref{prompts} of the Appendix and get more examples of how to extract/generate atomic facts from the sentence.

% This approach allows us to identify and assess hallucinations in terms of atomic levels and different categories. 
To achieve this, similar to the process of identifying descriptive sub-sentences, we also utilize the ChatGPT for the generation of atomic facts. This is because ChatGPT has shown a strong ability in information extraction~\cite{DBLP:journals/corr/abs-2302-10205}. More precisely, we annotate a set of $K_2$ examples for demonstrations and prompt the ChatGPT for atomic fact generation with $P'$ as follows:
% \begin{equation}
${E_i}= ChatGPT (A', P'), i \in [1,C]$, 
% \end{equation}
where $A'$ are all descriptive sentences identified in step 1, 
${E_i}=\{e_i^1,\cdots,e_i^{n_i}\}$ represents all $n_i$ atomic facts belonging to the $i$-th category, and $C$ stands for the total number of categories (\ie $C=5$ in our case) \textcolor[rgb]{0,0,0}{and the category include ``entity'', ``relation'', ``color'', ``count'', and ``others''}. 
% It is important to note that the set ${E_i}$ may be an empty set. For example, if an answer doesn't contain color information,  the corresponding set of the color category is empty.
% \xd{In what cases and why mentioning this?}
Further details regarding the specific prompt $P'$ utilized in this process can be found in Section~\ref{prompts} of the Appendix.

\paragraph{Fact Verification.}
In this stage, we compare the atomic facts derived above with the image to determine if the facts are faithful to the input visual information. Specifically, to calculate the \ours\ for the derived atomic facts, we first compute the score for each fact and then aggregate them to derive the overall score using the following formula:
\begin{equation}
\hat{s} = \frac{\sum_{i=1}^{C} \sum_{j=1}^{n_i} {w_i^j \cdot s (e_i^j, I )} }{ \sum_{i=1}^{C} {n_i} },
\end{equation}
where $\hat{s}$ represents the overall \ours\ of the answer $A$. The function $s(e_i^j, I)$ refers to the verification function (\ie Verifier), which measures whether $e_i^j$ can be supported by the input image $I$. The parameter $w_i^j$ is a weighted factor that can be used to assign different weights to different atomic facts for various tasks. 
To implement function $s(e_i^j, I)$, we resort to the Visual Entailment Model (VEM) (\eg OFA~\cite{ofa}), which is able to predict whether the image semantically entails the text. We elaborate on selections of the verifier models in Section \ref{43verifer}.
% Visual Entailment (VE) - is a task consisting of image-sentence pairs whereby a premise is defined by an image, rather than a natural language sentence as in traditional Textual Entailment tasks. The goal is to predict whether the image semantically entails the text.
In particular, when the output of the VEM is positive, indicating that the image $I$ semantically entails the text $e_i^j$ resulting in $s(e,I)=1$, and negative otherwise. 
% For example, [provide an example if needed]. 
In this work, we set all the weights $w_i^j$ to 1, following the setting of the existing work~\cite{factscore,DBLP:conf/eacl/KrishnaBKIDCL23}. 
% For a detailed overview of the overall method process, please refer to Algorithm~\ref{alg:overall}.
In addition, we further introduce a sentence-level \ours\ metric as follows, $\hat{s}_s = 1 - S_h/{S}$, 
% \begin{equation}
%     \hat{f}_s = 1 - \frac{C_h}{C}, 
% \end{equation}
where $S$ is the total number of descriptive sub-sentences in the answer and $S_h$ is the total number of descriptive sub-sentences with hallucinations.
